\documentclass[a4paper, 11pt]{article}
\usepackage[utf8]{inputenc}
\usepackage[T1]{fontenc}
\usepackage[italian]{babel}
\usepackage{geometry}
\usepackage{graphicx}
\usepackage{tabularx}
\usepackage{fancyhdr}
\usepackage{hyperref}
\usepackage[normalem]{ulem}

% Margini
\geometry{top=3cm, bottom=3cm, left=2.5cm, right=2.5cm}

% Intestazione e piè di pagina
\pagestyle{fancy}
\fancyhf{}
\fancyhead[L]{\small \textit{0x06tudio} -- Verbale di Riunione}
\fancyhead[R]{\small Anno Accademico 2026/2027}
\fancyfoot[C]{\thepage}
\renewcommand{\headrulewidth}{0.4pt}

\begin{document}
\thispagestyle{plain}

\newcommand{\issue}[1]{\href{https://github.com/0x06tudio/docs/issues/#1}{\textit{#1}}}
\newcommand{\linksott}[2]{\href{#1}{\uline{#2}}}
\newcommand{\verbale}{https://raw.githubusercontent.com/0x06tudio/docs/main/out/verbali/esterni/ve_2026_11_13.pdf}
\newcommand{\ddb}{https://raw.githubusercontent.com/0x06tudio/docs/main/out/diari_di_bordo/DiarioDiBordo5_ZeroSixStudio.pdf}


\begin{minipage}{0.6\textwidth}
    \vspace{-2cm}
    {\huge \textbf{0x06tudio}} \\[0.5em]
    {\large Progetto didattico di Ingegneria del Software} \\[0.3em]
    {\normalsize Università degli Studi di Padova}
\end{minipage}
\begin{minipage}{0.4\textwidth}
    \raggedleft
    \vspace{0pt}
    \includegraphics[width=6cm]{../../media/logo.png}
    \vspace{1cm}
    \small
    \href{mailto:0x06tudio@gmail.com}{0x06tudio@gmail.com}
\end{minipage}

\vspace{0.5cm}

\begin{center}
    {\LARGE \textbf{Verbale Riunione}}
\end{center}

\vspace{0.5cm}


\begin{tabularx}{\textwidth}{l X}
    \textbf{Versione:} & 0.1.0 \\
    \textbf{Stato:} & Bozza \\
    \textbf{Data:} & 2026/11/19 \\
    \textbf{Uso:} & Interno \\
    \textbf{Partecipanti (Presenti):} & Nicolò Beltrame, Pietro Brocadello, Catalin Chirilov, Davide Colabove, Andrea Maggio, Alberto Cavalletto \\
    \textbf{Destinatari:} & 0x06tudio \\
    \textbf{Modalità / Luogo:} & Online / Discord \\
    \textbf{Ora inizio:} & 15:30 \\
    \textbf{Ora fine:} & 17:30 \\
    \textbf{Durata totale:} & 2 ore \\
\end{tabularx}
\section*{Ordine del giorno}
\begin{enumerate}
    \item Ricerca e selezione del progetto di test per C9-SpecCheck\vspace{-0.2cm}
    \item Pianificazione delle attività tecniche in caso di approvazione\vspace{-0.2cm}
    \item Proseguimento della stesura dei documenti esterni
\end{enumerate}
\vspace{-0.5cm}

\section*{Registro delle modifiche}
\begin{tabularx}{\textwidth}{|c|c|X|c|c|}
    \hline
    \textbf{Versione} & \textbf{Data} & \textbf{Descrizione} & \textbf{Verificatore} & \textbf{Autore} \\ \hline
    0.1.0 & 2026/11/26 & Stesura del verbale  & / & Catalin Chirilov \\ \hline
\end{tabularx}

\newpage
\tableofcontents
\newpage


\section{Resoconto della discussione}

\subsection{Ricerca e selezione del progetto di test per C9-SpecCheck}
Il gruppo ha svolto un'attività di ricerca finalizzata all'individuazione di un progetto software già esistente da utilizzare come caso di studio per le attività di verifica previste dal nostro progetto (\textbf{C9-SpecCheck}). In particolare:
\begin{itemize}
    \item Sul \linksott{http://sarec.nd.edu/coest/datasets.html}{sito COEST} è stata individuata una raccolta di progetti utilizzabili come dataset di test.\vspace{-0.2cm}
    \item Tra i progetti disponibili nella raccolta, \textbf{eANCI} è emerso come il candidato più promettente.\vspace{-0.2cm}
    \item Il gruppo procederà ora a richiedere conferma e approvazione per l'utilizzo di \textbf{eANCI} sia al committente sia all'azienda proponente (Zucchetti S.p.A.).
\end{itemize}

\subsection{Pianificazione delle attività tecniche in caso di approvazione}
Subordinatamente alla conferma del progetto selezionato da parte di committente e proponente, il gruppo ha definito le seguenti fasi operative:
\begin{itemize}
    \item Implementazione di un flusso di conversione dei Casi d'Uso (UC) nei formati \textbf{Mermaid} e \textbf{XMI}, accompagnata dall'ampliamento delle relative specifiche dei requisiti.\vspace{-0.2cm}
    \item Una volta verificato e approvato il corretto funzionamento di questa fase preparatoria, avvio dello sviluppo delle funzionalità di \textit{parsing} e di verifica della corrispondenza tra casi d'uso, specifiche e codice.
\end{itemize}

\subsection{Proseguimento della stesura dei documenti esterni}
Parallelamente alle attività di esplorazione tecnica, il gruppo prosegue con la redazione dei seguenti documenti esterni:
\begin{itemize}
    \item \textit{Analisi dei Requisiti};\vspace{-0.2cm}
    \item \textit{Piano di Qualifica};\vspace{-0.2cm}
    \item \textit{Piano di Progetto};\vspace{-0.2cm}
    \item \textit{Norme di Progetto}.
\end{itemize}
\newpage


\section{Riepilogo delle decisioni e degli impegni (To-Do)}

\begin{tabularx}{\textwidth}{|c|X|X|}
    \hline
    \textbf{ID Decisione} & \textbf{Descrizione} & \textbf{Motivazione} \\ \hline
    D 9.1 & Caricamento della prima stesura del documento: \textit{Analisi dei Requisiti} & Archiviazione della documentazione approvata \\ \hline
    D 9.2 & Ricerca progetto svolta in gruppo & Migliorare la qualità della ricerca e possibilità di chiarimento più efficace \\ \hline
    D 9.3 & Attendere approvazione del proponente e del committente & Possibilità di ottenere feedback e chiarimenti\\ \hline
\end{tabularx}

\vspace{2cm}

\noindent
\begin{tabularx}{\textwidth}{|c|X|X|c|c|}
    \hline
    \textbf{Task ID} & \textbf{Descrizione} & \textbf{Assegnato a} & \textbf{Scadenza} & \textbf{Issue} \\ \hline
    TD 9.1 & Aggiornare Sprint S02 su sito web e nel documento "Piano di Progetto" & Pietro Brocadello & 2026/11/26 & \issue{80} \\ \hline
    TD 9.2 & Implementazione pipeline CI/CD per repository "SpecCheck" & Davide Colabove & 2026/12/18 & \issue{79} \\ \hline
    TD 9.3 & Redigere Diario di bordo 6 & Davide Colabove & 2026/11/29 & \issue{86} \\ \hline
    TD 9.4 & Proseguimento stesura documento "Piano di Qualifica" & Davide Colabove & 2026/12/18 & \issue{77} \\ \hline
    TD 9.5 & Proseguimento stesura documento "Analisi dei Requisiti" & Catalin Chirilov & 2026/12/18 & \issue{76} \\ \hline
    TD 9.6 & Proseguimento stesura documento "Norme di Progetto" & Andrea Maggio & 2026/12/18 & \issue{75} \\ \hline
    TD 9.7 & Proseguimento stesura documento "Piano di Progetto" & Pietro Broccadello & 2026/12/18 & \issue{74} \\ \hline
    TD 9.8 & Aggiungere visualizzazione "Studio tecnologie" nel sito web & Pietro Broccadello & 2026/11/29 & \issue{70} \\ \hline
    TD 9.9 & Invio mail al referente esterno Zucchetti & Tutti i membri & 2026/12/2 & \issue{87} \\ \hline
\end{tabularx}

\end{document}